\section{Conclusion}
% This study finds that the Texas shale gas boom had a positive impact on both household income and educational outcomes, particularly for low-income households. 
% Income gains were most concentrated among households in the bottom income quartile, and this group also experienced substantial increases in college enrollment and attainment.

% These findings challenge the traditional “Dutch Disease” narrative, which emphasizes the potential adverse effects of resource booms on long-term economic development. 
% Instead, our results suggest a potential “Dutch Blessing,” whereby resource-driven growth can generate positive spillovers for human capital accumulation—at least in the short to medium run.

% To ensure robustness, we employed both the staggered adoption difference-in-differences (SADiD) estimator and the Synthetic Control Method (SCM), leveraging their complementary strengths.

% Nonetheless, our analysis is subject to several limitations. 
% First, we focus primarily on short- to medium-term effects and do not account for longer-run dynamics. 
% Second, we do not model potential cross-state spillover effects, nor can we fully rule out bias from unobserved confounders.
% Finally, while we propose several mechanisms—such as increased public investment and improved household liquidity—we do not directly test them. 
% Future research could build on these results by incorporating firm-level data, migration patterns, and local government expenditure to better understand the channels through which resource booms affect economic and educational outcomes.

This study demonstrates that the shale gas boom in the United States had a positive impact on household income and educational outcomes, particularly for lower-income households, challenging the traditional Dutch Disease narrative. Using a staggered difference-in-differences framework and synthetic control methods, we find that the boom increased household income by approximately 4.8\% and 2.1\% 
for the first and second income quartiles, respectively, while also boosting college enrollment and high school attendance by 5\%p and 18\%p for the lowest quartile. 
These findings highlight a “Dutch Blessing,” where resource-driven economic growth enhanced both income and educational investments among low-income households, driven by increased employment opportunities, higher local tax revenues, and eased liquidity constraints. 
Although some pre-trend concerns and potential compositional changes warrant caution, robustness checks support the validity of our results. 
Future research could explore long-term impacts and further address potential biases from migration or other confounding factors to deepen our understanding of resource booms' heterogeneous effects.